% ============================================================
%  实验报告模板 — 王本熙 (25020007117)
%  编译方式: xelatex report.tex (运行两遍以解析交叉引用)
% ============================================================
\documentclass[a4paper,11pt]{ctexart}

% ---------- 页面布局 ----------
\usepackage[
  top=2.5cm, bottom=2.5cm, left=2.5cm, right=2.5cm,
  headheight=15pt, headsep=0.5cm, footskip=0.8cm
]{geometry}

% ---------- 数学与表格 ----------
\usepackage{amsmath, amssymb}
\usepackage{booktabs}
\usepackage{array}
\usepackage{multirow}

% ---------- 图形与链接 ----------
\usepackage{graphicx}
\usepackage[
  colorlinks=true, linkcolor=blue!70!black,
  citecolor=blue!70!black, urlcolor=teal!80!black
]{hyperref}
\usepackage{xcolor}

% ---------- 代码块 ----------
\usepackage{listings}
\definecolor{codebg}{RGB}{248,248,248}
\definecolor{codekw}{RGB}{0,0,180}
\definecolor{codecmt}{RGB}{0,128,0}
\definecolor{codestr}{RGB}{180,0,0}
\lstset{
  backgroundcolor=\color{codebg},
  basicstyle=\ttfamily\small,
  keywordstyle=\color{codekw}\bfseries,
  commentstyle=\color{codecmt},
  stringstyle=\color{codestr},
  numbers=left, numberstyle=\tiny\color{gray},
  frame=single, rulecolor=\color{gray!40},
  breaklines=true, showstringspaces=false,
  tabsize=2, columns=flexible,
  language=Python
}

% ---------- 页眉页脚 ----------
\usepackage{fancyhdr}
\pagestyle{fancy}
\fancyhf{}
\fancyhead[L]{\small 王本熙 \quad 25020007117}
\fancyhead[R]{\small 实验报告}
\fancyfoot[C]{\thepage}
\renewcommand{\headrulewidth}{0.4pt}

% ---------- 章节标题样式 ----------
\usepackage{titlesec}
\definecolor{secblue}{RGB}{30,80,160}
\titleformat{\section}
  {\Large\bfseries\color{secblue}}{第\thesection\,章}{0.6em}{}
\titleformat{\subsection}
  {\large\bfseries\color{secblue!80!black}}{\thesubsection}{0.5em}{}
\titleformat{\subsubsection}
  {\normalsize\bfseries\color{secblue!60!black}}{\thesubsubsection}{0.5em}{}

% ---------- 超链接边框样式微调 ----------
\hypersetup{
  pdftitle={实验报告},
  pdfauthor={王本熙}
}

% ============================================================
\begin{document}

% ---------- 封面 ----------
\begin{titlepage}
  \centering
  \vspace*{2cm}
  {\Huge\bfseries\color{secblue} 实验报告\par}
  \vspace{1cm}
  {\Large Tool Report\par}
  \vspace{3cm}
  \begin{tabular}{rl}
    \Large 姓\quad 名\,:\, & \Large 王本熙 \\[0.8em]
    \Large 学\quad 号\,:\, & \Large 25020007117 \\[0.8em]
    \Large 日\quad 期\,:\, & \Large \today \\[0.8em]
    \Large 平\quad 台\,:\, & \Large GitHub / Gitee \\
  \end{tabular}
  \vfill
  \rule{10cm}{0.4pt}
  \vspace*{1cm}
\end{titlepage}

% ---------- 目录 ----------
\tableofcontents
\newpage

% ============================================================
\section{实验目的}

本实验旨在掌握以下内容：
\begin{itemize}
  \item 使用 \LaTeX{} 编写结构化实验报告，包括公式、表格、图片与代码。
  \item 通过 \texttt{xelatex} 编译并生成 PDF，确保交叉引用不出现 \texttt{??}。
  \item 将代码与报告分批提交至 GitHub/Gitee，保留 commit 记录。
\end{itemize}

% ============================================================
\section{实验环境}

\begin{table}[htbp]
  \centering
  \caption{实验环境}\label{tab:env}
  \begin{tabular}{ll}
    \toprule
    项目 & 配置 \\
    \midrule
    操作系统 & Windows + WSL (Ubuntu) \\
    编辑器 & VS Code / Trae IDE \\
    编译器 & XeLaTeX \\
    版本管理 & Git \\
    远程仓库 & GitHub / Gitee \\
    \bottomrule
  \end{tabular}
\end{table}

% ============================================================
\section{实验结果}

% ----- 公式及引用 -----
本节展示带编号的公式与交叉引用。由式~\eqref{eq:pythagorean} 可得
$a^2 + b^2 = c^2$。

\begin{equation}
  a^2 + b^2 = c^2 \label{eq:pythagorean}
\end{equation}

% ----- 表格及引用 -----
表~\ref{tab:example} 展示了示例数据。

\begin{table}[htbp]
  \centering
  \caption{示例数据表}\label{tab:example}
  \begin{tabular}{cc}
    \toprule
    A & B \\
    \midrule
    1 & 2 \\
    3 & 4 \\
    \bottomrule
  \end{tabular}
\end{table}

% ============================================================
\section{代码与仓库}

项目已托管至 GitHub（请替换为你的实际地址）：
\url{https://github.com/USERNAME/damai-q04}

\subsection{Git 提交记录}

% 将截图放入本目录，取消下面两行注释
% \begin{figure}[htbp]
%   \centering
%   \includegraphics[width=0.9\textwidth]{commit-log.png}
%   \caption{Git commit 记录}
% \end{figure}

\subsection{仓库页面}

% \begin{figure}[htbp]
%   \centering
%   \includegraphics[width=0.9\textwidth]{repo-url.png}
%   \caption{GitHub 仓库页面}
% \end{figure}

\subsection{提交方式说明}

在 WSL (Ubuntu) 终端中执行分批提交，避免单次全量提交：

\begin{lstlisting}[language=bash]
# init repo
cd /mnt/d/damai_qiangpiao/q04
git init
git branch -M main
git remote add origin git@github.com:USERNAME/damai-q04.git

# commit in batches (no single bulk commit)
git add .gitignore  && git commit -m "chore: add gitignore"
git add report.tex  && git commit -m "docs: add report template"
git add report.pdf  && git commit -m "docs: add compiled pdf"
git push -u origin main

# show commit history
git log --oneline --graph --decorate
\end{lstlisting}

% ============================================================
\section{结论}

通过本实验，完成了基于 \LaTeX{} 的个性化实验报告模板设计，
公式与表格的交叉引用正常显示；同时通过 Git 分批提交
至 GitHub，保留了完整的 commit 记录。

\end{document}
